\documentclass[12pt]{article}
\usepackage[margin=1.15in]{geometry}
\usepackage{amsmath,amssymb,amsthm}
\usepackage{mathptmx}
\usepackage[expansion=false]{microtype}
\usepackage[T1]{fontenc}
\usepackage{setspace}
\usepackage[colorlinks=true,linkcolor=black,citecolor=black,urlcolor=black]{hyperref}
\setstretch{1.12}
\setlength{\parindent}{1.5em}
\setlength{\parskip}{0.2em}

\newtheorem{proposition}{Proposition}
\theoremstyle{definition}
\newtheorem{definition}{Definition}

\title{Authorization and Its Remainder}
\author{Flyxion\\ \small Independent Researcher}
\date{October 2026}

\begin{document}
\maketitle

\begin{abstract}
\noindent Jason Hickel's argument, as presented in a long-form interview with Miriam Fran\c{c}ois, locates the central failure of contemporary economies in the separation between productive capacity and collective authority over its use. This essay accepts that diagnosis as a serious one and then asks what it leaves unspecified. Democratizing the control of production answers the question of who may authorize an allocation. It does not answer who evaluates the result, who adjudicates among constituencies whose interests diverge, who may require correction after the fact, or what prevents the institution that acts from also certifying that it acted well. These functions form the remainder of authorization, and the part of the remainder that concerns judging results cannot legitimately inherit its standing from the decision it judges. The essay separates five institutional functions, states a simple result about self-certification, and examines the interview's treatments of workplace democracy, China, social democracy and degrowth in that light. It then extends the argument in three directions. First, drawing on Nietzsche's account of the contest and of ressentiment, it argues that the concentration of deference, and not only of ownership, provokes levelling responses at every scale, and that levelling destroys the evaluative function by removing the observers who differ from the group. Second, it argues that distance is a principal generator of the gap between those who decide and those who are affected, so that production organized locally, with information moving and matter staying, narrows the gap. Third, it argues that a society which removed the concept of garbage, by building objects that can be disassembled, that carry records of their composition, and that are produced in factories open to inspection, would provide evaluation with a channel to outcomes that the producer does not control. The essay does not claim that the interview's empirical assertions are false, and it does not establish that any alternative arrangement would perform better. Its claim is narrower: the argument as given does not support the confidence with which its institutional conclusion is stated, and several institutional features that would support it are identifiable.
\end{abstract}

\section{The claim the interview makes}

The interview \cite{hickelfrancois} is organized around a single contrast. On one side stands an extraordinary productive and technological capacity, described as greater than at any earlier point in history. On the other stands the persistence of conditions that, in the speaker's account, could be remedied quickly: people living on the streets of wealthy cities, inadequate public transport, insufficient renewable energy, and a large unmet demand for affordable housing. The explanation offered is not a shortage of knowledge or materials. It is that investment is directed by those who own investable assets and who allocate surplus according to expected profit, so that goods with high private returns are overproduced while goods with high social value but low private return are underproduced. Sport utility vehicles and fast fashion appear on one side of the ledger; public housing, durable products and renewable generation on the other.

From this diagnosis the interview draws a political conclusion that it states in several registers. Formal political democracy, in which citizens vote periodically, coexists with an economy in which decisions about what to produce, how to use labor and how to dispose of accumulated surplus are not democratic at all. A study of American policy outcomes published in 2014, evidently the work of Gilens and Page, is cited for the proposition that policy follows the preferences of affluent interests more closely than those of the majority \cite{gilens2014}. The analogy between consumer spending and voting, associated with Milton Friedman, is rejected on the ground that purchasing power is unequally distributed and that consumers cannot choose to have a public transit system built merely by being willing to pay for tickets. What is proposed instead is economic democracy: public control over finance and strategic industries, credit guidance of the kind used in mid-twentieth-century Europe, and worker ownership and governance of enterprises, with the Mondragon cooperatives and an Italian factory converted from automotive parts to bicycle components and solar panels offered as illustrations.

Around this core the interview arranges several further claims. Capitalism is historically specific, arising through enclosure and dispossession and expanding through colonial violence, and is therefore not a natural order. The Nordic social-democratic settlement depends on cheap labor and resources from the periphery and cannot be generalized. Digital platform power, which Varoufakis describes as techno-feudalism \cite{varoufakis2023}, is for the interview an outgrowth of the tendency of capital toward monopoly and not a post-capitalist form. Degrowth \cite{hickel2020} means reducing damaging and unnecessary production in rich countries while expanding necessary production, and it is supported by a maximum income taxed at one hundred percent and by an account of taxation as a tool for restraining excess demand in a sovereign currency issuer. China is described as an interesting case, with public control of finance and high reported satisfaction with government, though lacking worker control over production.

The aim of this essay is not to adjudicate these claims one by one. Several are empirical and contested, and some are stated in the interview with less support than their weight would require. The aim is to isolate a structural feature of the argument and to ask whether it survives scrutiny. That feature is the way authority over production is treated as the decisive variable.

\section{What the argument establishes}

It is worth stating plainly what is strong in the position before examining what is missing. Ownership of productive assets is not only an entitlement to income. It confers authority over which projects are funded, which are abandoned, who is employed, and where capacity is built. A government that votes for a housing program without access to land, finance, construction capacity or the political autonomy to command them has authorized something it cannot execute. The gap between formal authorization and effective control is real, and a theory of democracy that confines itself to legislatures and elections will miss it. The interview is right to press this point, and the observation that mid-century institutions for directing credit have been largely forgotten, even as technique, is a useful historical corrective.

The argument also correctly rejects the identification of profitability with social value. A profitable activity may impose costs that do not appear in the firm's accounts, may serve only those with sufficient purchasing power, and may crowd out activities that no private party can capture a return from. Public goods are the standard example. None of this requires a radical position; it is textbook welfare economics, and the interview's contribution is to insist on its scale.

What the argument establishes, then, is a thesis about the distribution of authorization: under existing arrangements the authority to allocate surplus is concentrated in a narrow class, and this concentration shapes which futures are materially available. What it does not establish is a thesis about the quality of allocations under a different distribution of authorization. The step from the first to the second is the step this essay examines.

\section{Five functions that economic democracy tends to run together}\label{sec:functions}

Consider any institution that directs productive capacity. Whether it is a bank, a cooperative, a ministry, a planning commission or an assembly, it performs, or delegates, some combination of five functions. It proposes: a candidate allocation is put forward. It authorizes: some body with standing approves the candidate. It executes: resources are committed and the project is carried out. It evaluates: the realized outcome is assessed against the purposes for which authorization was given. And it revises: the allocation, the rule that produced it, or the standing of those who produced it is altered in response to the evaluation. These can be written as maps. Let $x$ denote the state of the economy and the environment, $a$ an allocation, and $o$ the realized outcome. Then proposal selects $a$ from a feasible set $\mathcal{F}(x)$, authorization is a predicate $\alpha(a)$, execution produces $o$ from $a$ and exogenous conditions, evaluation is a report $r=v(o,a)$, and revision is a map from $r$ to a modified rule.

The interview concentrates its attention on authorization, and on who should hold it. Capital holds it now; the democratic alternative would distribute it across workers and citizens. This is a coherent and important proposal. It leaves the other four functions largely unspecified, and the omission is consequential because the functions can come apart. An allocation can be legitimately authorized and badly executed. It can be well executed and fail to serve its purpose because the purpose was misjudged. It can succeed on every internal measure while imposing costs on constituencies that did not take part in the authorizing procedure. And, most seriously, it can be evaluated only by the body that authorized and executed it.

The four functions other than authorization constitute what this essay calls its remainder. The remainder is not an implementation detail to be filled in after the deciding has been done, and its members do not stand in the same relation to the original decision. Execution may reasonably be delegated by authorization, since carrying out an approved plan is a task the approving body can assign. Evaluation and revision cannot be, because a body's authority to approve an allocation does not give it authority to certify its own result, to speak for persons it does not include, or to decide that it should be overruled. These functions require standing of their own, and an arrangement that lets them inherit standing from the authorization has not distributed power so much as concealed where it rests.

\begin{definition}
An evaluation is \emph{independent} of an authorizing body when the report $r$ depends on the realized outcome $o$ through channels the authorizing body does not control, so that $r$ carries information about $o$ beyond what is already determined by the authorized plan $a$.
\end{definition}

\begin{proposition}\label{prop:self}
If the report is a function of the authorized plan and the authorizing body's own records alone, so that $r=g(a)$, then the report is uninformative about the realized outcome conditional on the plan: $I(r;o\mid a)=0$.
\end{proposition}

The proof is immediate, since a function of $a$ carries no information about $o$ that is not already carried by $a$. The proposition is trivial, and its triviality is the point. It says that a body which certifies its own decisions using only its own records produces certifications that cannot distinguish a successful allocation from a failed one. This does not depend on bad faith. A well-intentioned body that measures what it has chosen to measure, in the units it chose, against targets it set, will find that its decisions were good, because the measurement apparatus was built to answer the question the decision already presupposed. Capital allocators exhibit this: a firm's accounts report profit, and profit is the criterion by which the firm's own investment decisions are justified. The interview identifies exactly this closure in the case of private investment. The difficulty is that the same closure is available to any institution that authorizes, executes and evaluates without separation, regardless of who holds the vote.

This is why the contrast between private and democratic control, taken alone, does not settle the matter. A cooperative that reports its own satisfaction, a public investment bank that reports its own lending performance, and a planning authority that reports its own plan fulfillment are in the same structural position as a firm that reports its own profit. Each may be better or worse than the others in what it selects for. None is thereby verified.

\section{The constituency problem}\label{sec:constituency}

The interview's examples of democratic production are drawn from workplaces. At Mondragon, workers elect managers, set pay ratios and share the yield of the enterprise. At the Italian factory, workers who had been dismissed reopened the plant and turned to producing goods they judged society needed. These examples are instructive, and they show that workers can govern firms and can redirect production toward purposes other than the owner's profit. They also illustrate a limitation of workplace democracy as a model for economic democracy generally.

Let an allocation $a$ affect a set $N(a)$ of persons, including present and future residents, consumers, neighbors, those downstream of any pollution, and those in other regions who supply inputs or absorb costs. Let $D(a)\subseteq N(a)$ be the set of persons with standing in the authorizing procedure. Workplace democracy sets $D(a)$ equal to the workers of the enterprise. The set $N(a)\setminus D(a)$ is the set of those affected without standing, and it is nonempty whenever the consequences of production extend beyond the workforce. A cooperative that votes to continue producing something harmful because its members depend on the income is procedurally democratic and may be socially harmful. This is not an objection peculiar to cooperatives; it is the ordinary structure of externalities, and it arises equally in a firm owned by shareholders. The point is that replacing shareholders with workers changes $D(a)$ without eliminating the gap.

The interview recognizes in passing that cooperatives operate under competitive pressures, and in the account of the Italian factory it presents workers choosing socially useful products. But the general case requires an answer to a question the interview does not pose directly: by what procedure are the claims of workers, consumers, neighbors and future persons to be reconciled when they conflict? Citizens' assemblies, which the interview cites favorably, are one candidate, and deliberative bodies may well move in the direction the interview describes. An assembly is, however, a forum in which a recommendation is recorded. A recommendation to expand public housing is not an allocation of land, financing, design and construction capacity, and the difference is exactly the difference between an inscription and a decision. The assembly's recommendation acquires binding force only through a further procedure, and the standing of that procedure is again a question of institutional design and not something inherited from the venue in which the preference was expressed.

None of this shows that democratic control of production would perform worse than the present arrangement. It shows that the claim that it would perform better rests on an assumption about how conflicts among constituencies are resolved, and that the interview supplies the assumption without the mechanism.

\section{Responsiveness, contestation and accountability}

The treatment of China exposes the issue most sharply. The interview presents survey evidence that large majorities of respondents in China report living in a democracy and report that the government serves their interests, and it describes feedback mechanisms through which development policy is said to respond to popular needs. It also acknowledges that worker control over production is lacking. When pressed about independent unions and the right to protest, the answer is that economic improvement is substantial, that the country remains a peripheral economy under global pressures, and that observers should remain open to seeing how events unfold.

The structural question is what the survey evidence is evidence of. It is evidence of reported satisfaction, and of reported perception that policy is responsive. These are properties of a system's relation to expressed preferences. Democracy as ordinarily understood includes at least two further properties: the capacity of citizens to contest decisions through channels the governing authority does not control, and mechanisms by which authorities can be sanctioned or replaced. Let $R$ denote responsiveness, $C$ the effective capacity to contest, and $A$ the existence of sanction and replacement. A system can score highly on $R$ while scoring poorly on $C$ and $A$, and the survey evidence by itself cannot distinguish a system with high $R$ and high $C$ from one with high $R$ and low $C$. The reason is the one stated in the previous section. A governing body that records preferences, interprets them, and reports on their satisfaction is in the position of the self-certifying evaluator. If respondents' answers are accurate reports of their satisfaction, which the interview argues on the basis of methodological studies, the evidence remains evidence about satisfaction. It does not by itself address whether the satisfaction would survive the availability of alternatives that the system does not permit to be organized.

The same analytical point applies in the other direction, and fairness requires saying so. The criticism the interview directs at Western liberal democracies, that they fail to translate majority preferences into policy, is a criticism that $R$ is low even where $C$ and $A$ are formally present. The two cases together show that the three properties vary independently, and that an account of economic democracy has to say which of them it is trying to secure and how it would know. Hickel's own position, as a democratic socialist who wants more worker control within the Chinese model, suggests that he treats $C$ and workplace voice as part of the goal. The interview does not, however, supply a criterion that distinguishes a responsive administration from a democratic one, and that criterion is what the comparison requires.

\section{The resentment of the exceptional}\label{sec:exceptional}

The interview locates illegitimate concentration of authority in a class defined by ownership, and the instruments it proposes address concentration of that kind: a maximum income, and ratios between the highest and lowest pay of the sort it describes at Mondragon. A second concentration is untouched by these instruments, and it is the concentration of deference. In any group, the person who is most capable, most persuasive, most generous with attention or most committed to a purpose can come to be listened to by others, and this happens without the ownership of anything. The question it raises for an account of answerable production is what a group does about such a person, and what the group's response does to the functions of Section~\ref{sec:functions}.

Nietzsche offers two accounts that bear on the question, and they should be kept apart. In ``Homer's Contest'' \cite{nietzsche1872} he reads the Greek practice of ostracism as an institution belonging to the logic of the contest. The Ephesians expelled Hermodorus, the best among them, on the principle that no one among them should be the best, and Nietzsche interprets the practice as a stimulus to renewed competition and not as a mere safety valve: a contest in which one competitor has overwhelmed the others has ceased to be a contest, and removing the preeminent individual reopens it. On this reading the levelling impulse has a function, which is to protect a plural order against the domination of a single figure. In the \emph{Genealogy of Morality} \cite{nietzsche1887} the same impulse appears as ressentiment, the reactive disposition of those who cannot act against what they resent and who respond by revaluing it as a fault. In \emph{Zarathustra} \cite{nietzsche1883} the great are pictured as surrounded by the buzzing of the small, who cannot forgive them for standing out, and the advice is withdrawal. The first account treats levelling as a safeguard with a cost, and the second treats it as a psychological formation with a cost. Both are relevant because the same behavior can be either, and from inside the group the two cannot be told apart.

The resentment does not track wealth. It arises in a poor neighborhood when everyone listens to one person, and the cause is structural. Listening is voluntary and revocable; authority is neither; and those who are not listened to experience the hardening of the first into the second as a loss of standing, even when the person listened to owns nothing. The resentment is therefore tracking something real, which is the conversion of deference into authority, the step from recorded preference to decision with standing that Section~\ref{sec:constituency} was concerned with. An egalitarian order that has capped wealth and wages but not the conversion of deference into authority has not addressed the portion of the problem that remains when ownership is equalized. It follows that ``the rich'' name one instance of a more general phenomenon. The richest member of a group, its cleverest, its most powerful and its most idealistic can each attract the same resentment, for the same structural reason, and the group's response to each is likely to be the same.

The difficulty is that the remedy by exclusion is indiscriminate. A group removes or silences the dominant and the merely unusual together, because it applies its norm and has no theory of which deviations threaten the contest. This bears directly on the evaluative function. Evaluation depends on observers who draw distinctions the authorizing body does not draw, and such observers are disproportionately those who differ from the group: the outsider, the dissenter, the holder of specialist or idiosyncratic knowledge, the idealist whose standard exceeds what the group is prepared to meet. If the group removes or silences observers whose distinctions are not shared by the majority, the set of distinctions that evaluation can draw contracts to the set the majority already holds. The report then becomes a function of what the authorizing body already believed, which is the condition of Proposition~\ref{prop:self} reached by a social path and not a bookkeeping one. Mill made the general point in \emph{On Liberty} \cite{mill1859}, arguing that social pressure toward conformity is a tyranny of its own and that the amount of eccentricity a society tolerates tracks the vigor of its members. The application to institutional design is that majority rule, extended to the question of who may evaluate, defeats evaluation.

Three things that the group's resentment fuses can be separated: being listened to, being obeyed, and being removable. An arrangement that keeps deference cheap to withdraw, authority costly to concentrate, and dissent costless to express allows a person to be listened to without becoming a bottleneck of the allocation. Positions of evaluation should be protected from the group they evaluate, through tenure, rotation, external appointment or selection by lot, so that an evaluator's standing does not depend on the approval of those evaluated. These protections create exceptional positions, which will attract the resentment the section describes. That is not a defect to be engineered away but a standing cost, and a society that cannot bear it will lose its evaluators. The maximum income and the pay ratios the interview endorses address material differences, and nothing here argues against them. The point is only that they leave the structure of deference as they found it.

\section{Profit as a signal, and what replaces it}

The interview is emphatic that profit is a poor measure of social value, and as a measure of social value it is. But a measure can be poor for one purpose and informative for another. Profit and loss register, imperfectly and with distortion from market power, subsidy and unequal purchasing power, whether an organization is consuming more than others are willing to give it in exchange. That is a signal about opportunity cost. Rejecting profit as the objective of investment does not remove the need for some means of comparing alternative uses of scarce resources.

The difficulty can be made concrete. A community with a given stock of skilled labor and materials may be able to build either a hospital or several hundred homes. Both are socially necessary. Democratic authorization of the objective of meeting needs does not determine the allocation, and neither does the statement that profit is not the criterion. Some procedure of valuation under scarcity is required, and it must be able to represent trade-offs, uncertainty and the information held by those distributed across the economy. This is a version of the problem to which Hayek drew attention in discussing the dispersal of knowledge \cite{hayek1945}. The relevant point is not that markets solve it and planning cannot, a conclusion that does not follow, and not that democratic bodies are incapable of building valuation procedures, which they plainly can. It is that the interview rejects one coordinating device and describes the values to be served, without describing the replacement device. A rejection of profit as an objective is stronger than an account of what performs the evaluative work profit does imperfectly.

The accounting problem runs deeper than the choice of objective. Any ledger that rewards a recorded quantity can diverge from the purpose that quantity was meant to track. Gross output rises when a building is constructed, when it requires repair, and when it is demolished and replaced, and none of these entries establishes that the built environment has improved. The interview's critique of aggregate growth as a measure of welfare is an instance of this divergence. What it implies is that any proposed replacement ledger must be examined for the same property, including ledgers kept by public enterprises and cooperatives. A metric of housing units completed can rise while the quality, location and durability of housing falls. A metric chosen democratically is still a metric, and the relevant question for any metric is what would show it to be diverging from the purpose it stands for.

\section{Foreclosure and reversibility}\label{sec:foreclosure}

A further consideration concerns the direction of time. Decisions about production create future constraints. An infrastructure built around a particular energy source, a technology that binds users to one supplier, a pattern of land use that removes land from future public use: each converts a present choice into a limit on later choices. Let $\mathcal{A}_{t+1}(a_t)$ denote the set of actions that remain available after allocation $a_t$. Two allocations that deliver similar immediate benefits may leave very different sets of future options. An allocation procedure that attends only to present benefit, whether measured by profit or by democratically expressed preference, can shrink this set without any party intending it.

This is relevant to the interview's ecological argument, which is its most forceful. The case for reducing damaging production is in part a case about irreversible loss, and it is a strong one. But the same consideration applies to proposed remedies. A transition that is legitimately authorized and later found mistaken should remain correctable, which requires that the decision not remove the capacity to undo it and that some body other than its authors have standing to say that it should be undone. A society that governs its future does not only choose who controls production at a given moment. It preserves the conditions under which production decisions can be challenged and their consequences assessed. That formulation is not a rival to the interview's view; it is a requirement any version of economic democracy would need to meet to be taken seriously as self-government over time, and it is one that private control fails to meet in specific ways and that democratic control could fail to meet in others.

\section{Degrowth, classification and transition}

The interview's account of degrowth is more careful than its popular reception. It rejects the idea that poorer countries should be restrained, locates the argument in the high-consumption economies, and distinguishes the contraction of damaging and unnecessary production from the expansion of housing, renewable energy, public services, public transport and food provision. As stated, this is a theory of selective reallocation and not of general contraction, and it should be evaluated as that.

Its principal difficulty is classificatory. Whether an activity is necessary or unnecessary depends on context. A large vehicle may be an urban extravagance or a practical necessity in a remote community. A data center may serve manipulative advertising or essential research. A large house may be a luxury residence or a multigenerational household. The environmental consequence of production depends on technology, geography, durability, supply chains and substitution. Selective contraction therefore requires, in addition to a political decision to contract, an institutional means of assessing the marginal social and ecological consequences of particular activities, and a means of managing what happens to workers and communities whose industries shrink. The interview supplies a rationale for reallocation and an indication that citizens' assemblies tend to favor it. It does not supply the evaluative and transitional machinery. This is not surprising in a conversational format, but it matters, because the classification step is exactly where the five functions of Section~\ref{sec:functions} come into play. Someone proposes the classification, some body authorizes it, some agency executes it, and the question of who assesses whether the resulting reallocation achieved what it was meant to achieve returns in full.

A similar remark applies to the proposal for a maximum income. The concern about unlimited accumulation as unlimited command over labor and resources is intelligible, and the account of money as claim on real resources has force. The instrument is stated at a level of generality that leaves open the treatment of capital gains, wealth, business ownership, non-cash compensation and international mobility, and the interview refers readers to book-length treatments for these. The same is true of the monetary argument that a sovereign issuer is not financially constrained in the manner of a household. That argument is a useful corrective to household analogies, but monetary sovereignty does not remove real resource constraints, import dependence or exchange-rate pressure, and the interview acknowledges some of these in its discussion of capital flight.

\section{Unequal exchange and the limits of generalization}

The argument about the Nordic countries makes a different kind of claim, and it deserves separate treatment because it is offered as a reason that social democracy cannot be a universal solution. The claim is that domestic prosperity and concessions to labor in high-income countries are made possible by access to cheap labor and resources elsewhere, and that violence against organizers and protesters in peripheral economies is the counterpart of the apparent calm of the core. The interview cites empirical studies of unequal exchange and does not describe their methods.

The question it raises is one of counterfactual attribution. How much of high-income countries' welfare provision depends on unequal exchange, and how much on domestic productivity, institutions, taxation and accumulated capital? The existence of unequal exchange does not by itself show that generous provision could not be generalized under more equal terms of trade. These questions are open, and answering them requires the kind of methodological detail a conversation cannot carry. The claim should therefore be read as a hypothesis with serious supporting literature and not as a settled result. The same caution applies to the historical narrative that presents capitalism as a reaction to the success of late-medieval peasant revolts: it is an intelligible interpretation of a heterogeneous set of developments, and compressing it into one causal story gains rhetorical force at the cost of precision.

The analytical contribution of the unequal-exchange argument for present purposes is different from its empirical status. It asserts that the boundary of the constituency matters: an arrangement that looks answerable within one polity may be unanswerable to those outside it who bear its costs. This is the constituency problem of Section~\ref{sec:constituency} at global scale. It strengthens the case that standing in the authorizing procedure and the set of persons affected can diverge, and it suggests that no domestic democratization, however complete, resolves the gap by itself.

\section{Locality and the narrowing of the constituency gap}\label{sec:locality}

Distance is one of the principal generators of the set $N(a)\setminus D(a)$. A decision about production made in one place imposes costs and benefits on persons elsewhere: extraction occurs in one region, assembly in a second, use in a third, disposal in a fourth. The interview's account of unequal exchange is in effect the observation that this separation is arranged so that costs fall on those furthest from the authorizing body. Reducing the separation reduces the gap. If production is carried out at a scale and in a place where those who use it, work in it and live beside it are largely the same persons, then $D(a)$ approaches $N(a)$ without any change in voting rules, and the outcomes of the allocation are observable by those with standing to judge them. Proximity is thus an institutional variable and not merely a logistical one, since it determines how much of the remainder can be performed by the people who bear the consequences.

The limit case of this principle can be stated as a rule that information moves and matter stays. Designs, instructions and fabrication procedures circulate widely, while physical goods are made near where they are used and food is grown near where it is eaten. In the limit, a community fabricates the objects it uses, including electronic ones, and grows its own food, and the freight that now connects distant stages of production largely disappears. The essay does not claim that this limit can be reached. Leading-edge semiconductor fabrication is among the most complex industrial installations in existence and depends on specialized materials and a deep supply chain. Climate, water and soil constrain which food can be grown where. Some minerals occur only in particular places, and economies of scale are real in some processes. The claim is the weaker one that movement toward the limit is movement toward a smaller constituency gap, so that the costs of distance should be counted in the assessment of any production arrangement alongside the costs of local production, and that the balance should be assessed and not assumed in either direction.

The interview is introduced with the observation that consumer flights continue to climb even as concern about ecological collapse grows. A hypothesis, which this essay offers as a hypothesis and not as a finding, is that part of long-distance leisure travel is flight from a local environment experienced as finished, closed and mundane, and that borders sharpen the experience by defining the elsewhere as the place where different things are permitted and different objects are made. Travel has many motives, among them family, work, curiosity and climate, and nothing here ranks them. If the hypothesis holds in part, however, an environment whose making is open to inspection and variation reduces one motive for departure at the same time as it reduces freight, and the two reductions are connected through the same feature: an economy whose production is local and legible is one in which the mundane is less opaque.

\section{Objects that keep their own accounts}\label{sec:objects}

Garbage is an accounting category before it is a physical one. An object becomes garbage at the point where it leaves every ledger: the producer's account closed at sale, the owner's closed at disposal, and no one's account open thereafter. The costs of that exit are the externalities discussed above, and they accrue at the boundary where the ledger stops. A society that removed the concept of garbage would not be one in which nothing is discarded, which no process can achieve, since reuse and recycling dissipate material and energy. It would be one in which no object leaves the ledger: every manufactured thing is held in the account of a custodian from fabrication until its materials are returned to use, every residue is a named stock with a named destination, and unaccounted dissipation is reported as a loss and not as nothing. This is the reversibility condition of Section~\ref{sec:foreclosure} applied at the scale of objects. A thing that can be taken apart leaves open the actions of repair, reuse and repurposing, which a sealed thing forecloses, so that the set of admissible actions after fabrication is larger.

Three requirements follow. Each manufactured object can be disassembled with the means available to those who receive it. Each is scanned and labeled at fabrication with a record of its materials, components, assembly sequence and provenance. And each is returned to a custodian when it no longer serves. Only the second bears directly on evaluation, and it must be stated carefully, because it is easily misread. A label written by the producer is the producer's claim, and a society of labels is a society of self-reports, which Proposition~\ref{prop:self} says are uninformative. The value of the record lies in the fact that it can be tested against the object by parties other than the one who wrote it. A device labeled as repairable that cannot be opened, or labeled as containing one material and found to contain another, produces a discrepancy that anyone with the means to open it can detect. The object is the evidence and the label is the claim, and independent evaluation is the comparison of the two by someone the producer does not control. In the vocabulary of Definition~1, the object gives the report $r$ a channel to the realized outcome $o$ that does not pass through the authorizing body's records. The interview complains that products are short-lived and that durable products are underproduced because they are not profitable. The three requirements are a candidate institutional counterpart, and the record is what allows a claim of durability to be checked and not merely asserted. Records of this kind supply data to valuation without replacing it, since they do not settle how a hospital is to be compared with several hundred homes.

A proposal for embedding records in the material of the object itself bears on this requirement. Laser-etched holographic patterns written into stone, brick, wood, cloth or the structure of a printed object, with many adjustable parameters governing how much information they hold, would allow the record to travel with the object and to remain readable after the object has been separated from the paperwork, the database and the firm that made it. If records can be written into ordinary materials at fabrication, the account of an object need not depend on a registry that someone could close, rewrite or lose, and disassembly could proceed from the object's own instructions. The capacity and robustness of such encodings are matters for engineering analysis and remain open. The essay uses the proposal only to illustrate a principle, which is that the evidence and its record can be located together.

The second element is the legibility of production. A navigation system developed for exploring a large body of texts as a three-dimensional structure, modeled on the cutaway automap of the game \emph{Descent}, indicates the form such legibility can take: movement through a structure while its layers remain visible and the history of the path is kept. If schools conduct tours through the working interiors of the factories that supply them, students, apprentices and neighbors acquire standing to observe production. This widens the set of persons able to evaluate and trains the observers that the evaluative function requires, and it addresses the constituency gap of Section~\ref{sec:constituency} from the side of observation, since persons who can see how a thing is made are better placed to say whether it should be. Two cautions apply. A tour offered by the producer is a presentation chosen by the producer, and it shares the weakness of any self-report unless visitors can go where the producer did not choose and can compare what they see with what is made. A tour combined with the disassembly of the product is stronger than either alone. And some processes are hazardous, so that access requires physical separation, viewing galleries and recorded or instrumented views. These are engineering questions that do not alter the principle, and they should be settled without making the exceptions larger than the rule.

\section{Why the distinction matters}

It may be objected that the demands of this essay are unreasonable, since no institution can satisfy independent evaluation, constituency reconciliation, reversibility and valuation under scarcity, and that insisting on them amounts to a defense of the status quo. That objection has the weight of a real point: perfection is not the standard, and the present system fails on every one of these dimensions as well. But the essay's claim is comparative and structural. It does not say that economic democracy should be rejected or that profit-driven investment is adequate. It says that the case for the former, as presented, places its entire weight on a single variable, and that the variable by itself does not distinguish institutions that are answerable to the consequences of their decisions from those that are not.

The historical examples the interview invokes are consistent with this reading. The improvements it credits to organized movements, namely labor rights, public services, sanitation, the weekend and the franchise, were not only transfers of authority. They were also the creation of independent mechanisms through which claims could be pressed and outcomes checked: inspectorates, courts, free unions, independent press, public statistics, and the protected standing of the people who did the checking. Authority was redistributed, and standing to contest and to evaluate was extended with it. A program of economic democracy that attended only to the first of these would reproduce half of what made the earlier gains durable.

There is also a point about the relation between this essay and the interview that should be stated openly. The interview is a conversation of limited length, conducted for a general audience, and it moves quickly among historical interpretation, empirical report, institutional proposal and normative commitment, often without marking the transitions. The omissions identified here may be reasonable omissions in that format, and Hickel's written work may address several of them. The essay's concern is with what the argument, as given, supports, and whether the claim that democratization of production is the solution can be separated from the claim that it is a necessary part of a solution. The latter is considerably more defensible than the former.

\section{Conclusion}

The interview's most durable contribution is the identification of a gap between possessing productive capacity and possessing legitimate authority to direct it. That gap is real, and its description is a better account of contemporary institutional failure than many that confine themselves to the language of markets and states. The remedy it proposes, a wider distribution of authorization over finance and production, addresses one of the five functions through which any institution directs resources. It does not address the remainder, and the remainder is where democratic economic organization will succeed or fail in practice. Evaluation that is independent of authorization requires observers who are protected from the groups they observe, because a society that levels its exceptional members, whether the richest, the cleverest or the most idealistic, removes the people who draw the distinctions evaluation depends on. Procedures that reconcile the claims of those affected with those who vote are easier to build where production is near its consequences, so that the set of affected persons and the set with standing largely coincide. And the evidence on which any evaluation rests is stronger where objects can be opened, where they carry records that can be tested against their contents, and where the making of things can be seen. Without these conditions, a democratically controlled economy is exposed to the failure the interview diagnoses in the present one: the selector's account of its own success is accepted as the measure of success. The question that remains after the interview is therefore not only who should hold authority over production, but what arrangement makes any holder of that authority answerable to the future it produces.

\begin{thebibliography}{9}
\bibitem{gilens2014} M.~Gilens and B.~I.~Page, ``Testing Theories of American Politics: Elites, Interest Groups, and Average Citizens,'' \emph{Perspectives on Politics} 12(3), 564--581, 2014.
\bibitem{hayek1945} F.~A.~Hayek, ``The Use of Knowledge in Society,'' \emph{American Economic Review} 35(4), 519--530, 1945.
\bibitem{hickel2020} J.~Hickel, \emph{Less Is More: How Degrowth Will Save the World}, Windmill Books, 2020.
\bibitem{hickelfrancois} J.~Hickel in conversation with M.~Fran\c{c}ois, \emph{The Tea}, interview on capitalism, democracy and economic democracy, 2026.
\bibitem{mill1859} J.~S.~Mill, \emph{On Liberty}, 1859, chapter 3, ``Of Individuality, as One of the Elements of Well-Being.''
\bibitem{nietzsche1872} F.~Nietzsche, ``Homer's Contest'' (1872), in \emph{The Portable Nietzsche}, trans.\ W.~Kaufmann, Viking, 1954.
\bibitem{nietzsche1887} F.~Nietzsche, \emph{On the Genealogy of Morality} (1887), first essay.
\bibitem{nietzsche1883} F.~Nietzsche, \emph{Thus Spoke Zarathustra} (1883--85), part I, ``On the Flies of the Marketplace.''
\bibitem{varoufakis2023} Y.~Varoufakis, \emph{Technofeudalism: What Killed Capitalism}, Bodley Head, 2023.
\end{thebibliography}

\end{document}

